\documentclass[12pt]{ximera}
\title{Exam 2}
\begin{document}
\begin{abstract}
Exam 2, covering Sections 7.5--8.2: what the multiplication principle can and cannot
count, testing independence from survey data, four counting techniques side by side,
and Bayes' Theorem for rain and wind.
\end{abstract}
\maketitle

\section*{Exam 2}

\begin{problem}
Assess whether each of the following claims is true or false.

\begin{prompt}
\textbf{(a)} Claim: the multiplication principle can be used to count any situation that
can be framed as a series of choices made one at a time.

\begin{multipleChoice}
\choice{True}
\choice[correct]{False}
\end{multipleChoice}

\begin{feedback}
False. Multiplying stage counts counts \emph{ordered} sequences of choices. When order
does not matter, that overcounts: choosing a committee of $3$ from $20$ can be framed as
three one-at-a-time choices, but $20\cdot 19\cdot 18=6840$ counts each committee $3!=6$
times. The correct count is $C(20,3)=1140$.
\end{feedback}
\end{prompt}

\begin{prompt}
\textbf{(b)} Claim: the multiplication principle does NOT work in situations where order
matters.

\begin{multipleChoice}
\choice{True}
\choice[correct]{False}
\end{multipleChoice}

\begin{feedback}
False --- it is the reverse. The multiplication principle is at its most natural when
order matters; the permutation formula $P(n,r)=n(n-1)\cdots(n-r+1)$ is a direct
application of it.
\end{feedback}
\end{prompt}
\end{problem}

\begin{problem}
Suppose that in a survey of $64$ students, $36$ report they like turtles and $48$ report
they like hedgehogs; $27$ of the students report they like \emph{both}.

\begin{prompt}
\textbf{(a)} Among the students who like turtles, what is the probability that they also
like hedgehogs? Enter a fraction in lowest terms.

$\answer{3/4}$

\begin{feedback}
Of the $36$ turtle-likers, $27$ also like hedgehogs: $P(H\mid T)=\frac{27}{36}=\frac34$.
\end{feedback}
\end{prompt}

\begin{prompt}
\textbf{(b)} Two events are independent if assuming one of them has no impact on the
probability of the other: $P(E\mid F)=P(E)$. In this survey, are liking turtles and
liking hedgehogs independent events?

\begin{multipleChoice}
\choice[correct]{Yes}
\choice{No}
\end{multipleChoice}

\begin{feedback}
Yes. Overall, $P(H)=\frac{48}{64}=\frac34$, which is exactly $P(H\mid T)$ from part (a).
Knowing a student likes turtles does not change the chance they like hedgehogs.

Equivalently, $P(T)\cdot P(H)=\frac{36}{64}\cdot\frac{48}{64}=\frac{27}{64}=P(T\cap H)$.
\end{feedback}
\end{prompt}
\end{problem}

\begin{problem}
Use counting techniques to answer the following. For each, identify whether you're using
the multiplication principle, permutations, \emph{distinguishable} permutations, or
combinations.

\begin{prompt}
\textbf{(a)} How many different starting lineups of five basketball players from a team
of $15$ are possible, if we consider a lineup just a subset?

\begin{multipleChoice}
\choice{Multiplication principle}
\choice{Permutations}
\choice{Distinguishable permutations}
\choice[correct]{Combinations}
\end{multipleChoice}

The count is $\answer{3003}$.

\begin{feedback}
A subset of $5$ from $15$: $C(15,5)=3003$.
\end{feedback}
\end{prompt}

\begin{prompt}
\textbf{(b)} How many different rankings are possible if you select and rank your three
favorite songs from an album of $20$ different songs?

\begin{multipleChoice}
\choice{Multiplication principle}
\choice[correct]{Permutations}
\choice{Distinguishable permutations}
\choice{Combinations}
\end{multipleChoice}

The count is $\answer{6840}$.

\begin{feedback}
A ranking is ordered with no repeats: $P(20,3)=20\cdot 19\cdot 18=6840$.
\end{feedback}
\end{prompt}

\begin{prompt}
\textbf{(c)} How many distinct patterns can you make by arranging $7$ yellow, $4$ red,
and $5$ blue painted tiles all in a row?

\begin{multipleChoice}
\choice{Multiplication principle}
\choice{Permutations}
\choice[correct]{Distinguishable permutations}
\choice{Combinations}
\end{multipleChoice}

The count is $\answer{1441440}$.

\begin{hint}
There are $16$ tiles, but tiles of the same color are interchangeable. Divide out the
rearrangements within each color.
\end{hint}

\begin{feedback}
\[
\frac{16!}{7!\,4!\,5!}=1{,}441{,}440.
\]
The $16!$ counts arrangements as if every tile were different; dividing by $7!$, $4!$,
and $5!$ removes the swaps of same-colored tiles, which do not change the pattern.
\end{feedback}
\end{prompt}

\begin{prompt}
\textbf{(d)} How many different $6$-digit passcodes are possible if each digit is a
number from $0$ to $9$?

\begin{multipleChoice}
\choice[correct]{Multiplication principle}
\choice{Permutations}
\choice{Distinguishable permutations}
\choice{Combinations}
\end{multipleChoice}

The count is $\answer{1000000}$.

\begin{feedback}
Each of the $6$ digits is chosen independently from $10$, with repeats allowed:
$10^6=1{,}000{,}000$.
\end{feedback}
\end{prompt}
\end{problem}

\begin{problem}
You've been inside all day, but you know that the chance of high winds today is $40\%$.
The chance of rain if there are high winds is $40\%$, and the chance of rain if there are
\emph{no} high winds is $30\%$. Your friend comes over and mentions that it is raining.
Given this, what is the likelihood that there are high winds today?

\begin{prompt}
\textbf{(a)} Build the tree diagram. Which event belongs on the first set of branches?

\begin{multipleChoice}
\choice[correct]{Wind first, then rain}
\choice{Rain first, then wind}
\end{multipleChoice}

Branch probabilities:

$P(W)=\answer[tolerance=0.001]{0.4}$ \qquad $P(\olsi{W})=\answer[tolerance=0.001]{0.6}$

\smallskip
$P(R\mid W)=\answer[tolerance=0.001]{0.4}$ \qquad $P(\olsi{R}\mid W)=\answer[tolerance=0.001]{0.6}$

\smallskip
$P(R\mid \olsi{W})=\answer[tolerance=0.001]{0.3}$ \qquad
$P(\olsi{R}\mid \olsi{W})=\answer[tolerance=0.001]{0.7}$

\begin{feedback}
The information given is $P(W)$ and then the chance of rain \emph{given} the wind
situation, so wind goes on the first branches and rain on the second. Each pair of
branches from a point adds to $1$.
\end{feedback}
\end{prompt}

\begin{prompt}
\textbf{(b)} Multiply along each path.

$P(W\cap R)=\answer[tolerance=0.001]{0.16}$ \qquad
$P(W\cap \olsi{R})=\answer[tolerance=0.001]{0.24}$

\smallskip
$P(\olsi{W}\cap R)=\answer[tolerance=0.001]{0.18}$ \qquad
$P(\olsi{W}\cap \olsi{R})=\answer[tolerance=0.001]{0.42}$

\begin{feedback}
$(0.4)(0.4)=0.16$, $(0.4)(0.6)=0.24$, $(0.6)(0.3)=0.18$, $(0.6)(0.7)=0.42$, which total
$1$. \checkmark
\end{feedback}
\end{prompt}

\begin{prompt}
\textbf{(c)} Knowing that it \emph{is} raining, what is the probability that there are
high winds today? Give a percentage rounded to two decimal places.

$\answer[tolerance=0.005]{47.06}\%$

\begin{feedback}
Two paths end in rain, so $P(R)=0.16+0.18=0.34$, and
\[
P(W\mid R)=\frac{P(W\cap R)}{P(R)}=\frac{0.16}{0.34}=\frac{8}{17}\approx 47.06\%.
\]
Learning that it is raining raises the chance of high winds from $40\%$ to about $47\%$,
because rain is somewhat more likely when it is windy.
\end{feedback}
\end{prompt}
\end{problem}

\end{document}
